\input{../header}
\hyphenpenalty=5000

\rhead{Your name: \rule{8cm}{0.15mm}}

\renewcommand{\thefootnote}{\fnsymbol{footnote}}

\begin{document}

\everymath{\displaystyle}

\section{Learning target D4, version 1}

Find general solutions to separable differential equations, and use initial values to determine the value of the constant of integration.

\hrulefill

Suppose that a cylindrical water tank with a hole in the bottom is filled with water. The water, of course, will leak out and the height of the water will decrease. Let $h(t)$ denote the height of the water. A physical principle called \textit{Torricelli's Law} implies that the height decreases at a rate proportional to the square root of the height:
\footnote{Conservation of energy: the kinetic energy of water leaving the tank is equal to the change in GPE.}

\[\dfrac{dh}{dt} = -k \sqrt{h}, \quad k > 0.
\]

\begin{enumerate}
	\item Suppose you have a tank for which the height decreases at a rate of 20 inches per minute when the water is filled to a depth of 100 inches. Find the value of $k$.
	\vfill
	\item Now solve the differential equation.
	\vfill
	\vfill
	\item Finally, use the initial condition $h(0) = 100$ to find the value of the constant of integration.
	\vfill
\end{enumerate}
%%%%%%%%%%%%%%%%%%%%%%%%%%%%%%%%%%%%%%%%%%%%%%%%%%%%%%%
\pagebreak
\section{Learning target D4, version 2}

Find general solutions to separable differential equations, and use initial values to determine the value of the constant of integration.

\hrulefill

Say you take a soda out of the refrigerator, where it is 38$^\circ$F, and you put it on the counter in a 73$^\circ$F room. According to Newton's Law of Cooling (Slash Heating), the soda's temperature $T$ changes at a rate proportional to the difference between the soda's temperature and room temperature:

\[\dfrac{dT}{dt} = -k \cdot(T-73), \quad k > 0.
\]

\begin{enumerate}
	\item What's the initial condition $T(0)$?
	\vspace{2em}
	\item Solve the differential equation, and find the value of the constant of integration.
	\vfill
	\vfill
	\item Five minutes after you take the soda out of the fridge, you measure that its temperature is 41$^\circ$F.\footnote{I have no idea if this is realistic.} Use this extra information to figure out the value of $k$.
	\vfill
\end{enumerate}
%%%%%%%%%%%%%%%%%%%%%%%%%%%%%%%%%%%%%%%%%%%%%%%%%%%%%%%
\pagebreak
\section{Learning target D4, version 3}

Find general solutions to separable differential equations, and use initial values to determine the value of the constant of integration.

\hrulefill

The Gompertz equation is a model that is used to describe the growth of certain populations. Suppose that $P(t)$ is the population of some organism, measured in thousands, and that

\[\dfrac{dP}{dt} = -P \cdot \ln\left(\tfrac{1}{3}P\right).
\]

\begin{enumerate}
	\item Solve this differential equation. (Hint: try $u = \ln(\tfrac{1}{3}P)$)

	\vfill

	\item If $P(0) = 1$, what's the value of the constant of integration? What happens to the population long-term?
	\vfill
	\item What if $P(0) = 6$?
	\vfill
	\item (Bonus!) Look at a slope field.
\end{enumerate}
%%%%%%%%%%%%%%%%%%%%%%%%%%%%%%%%%%%%%%%%%%%%%%%%%%%%%%%
\pagebreak
\section{Learning target D4, version 4}

Find general solutions to separable differential equations, and use initial values to determine the value of the constant of integration.

\hrulefill

When ice forms on a lake, the water on the surface freezes first, and then the layer of ice insulates the water below it. Therefore, the thickness of the ice $I(t)$ increases at a rate that's inversely proportional to how thick it currently is:
\[\dfrac{dI}{dt} = \dfrac{k}{I}, \text{ for some } k > 0.\]

\begin{enumerate}
	\item Solve this differential equation. (Hint: notice that $I(t)$ is always non-negative.)

	\vfill

	\item Measure $t$ from the moment the ice started freezing. What then is $I(0)$? Use this to solve for the constant of integration.
	\vfill
	\item Use Desmos to investigate the effect of a couple of different values of $k$ on the graph of $I(t)$. What real-world factors do you think contribute to the value of $k$?
\end{enumerate}
%%%%%%%%%%%%%%%%%%%%%%%%%%%%%%%%%%%%%%%%%%%%%%%%%%%%%%%

\end{document}